\SourceLeaf{35}
\fancyhead[L]{\small Principia generalia theoriae figurae fluidorum}
\renewcommand{\Article}[1]{\par\medskip\Needspace{4\baselineskip}\pdfbookmark[1]{Art. #1}{fluidorum-art-#1}\centerline{\textbf{#1.}}\nobreak\smallskip}
\pdfbookmark[0]{Principia generalia theoriae figurae fluidorum in statu aequilibrii}{fluidorum-body}
\par\medskip
\centerline{PRINCIPIA GENERALIA}
\centerline{THEORIAE FIGURAE FLUIDORUM}
\centerline{IN STATU AEQUILIBRII.}
\medskip
\expandafter\def\csname HistoricalNoteText4\endcsname{Constat, maximam attractionem, quam massa homogenea data in punctum datum secundum illam legem exercere potest, esse ad attractionem, quam eadem massa in figuram sphaericam redacta exercet in punctum in superficie positum, ut $3$ ad $\BareRoot[3]{25}$: posterior vero attractio cum gravitate facile comparatur.}
\expandafter\def\csname HistoricalNoteText5\endcsname{Ita iudicandum de plerisque obloquutionibus in ephemeridibus Ticinensibus (Giornale di fisica etc. T. 9) prolatis, quibus scite respondit clar. \textsc{Petit} in Annales de chimie et de physique T. 4.}
Vires ascensionem vel depressionem fluidorum in tubis capillaribus gubernantes primus acute et accurate enumeravit sagax \textsc{Clairaut}, sed quum legem virium omnino intactam liquerit, nihil fructus ad explicationem mathematicam phaenomenorum ex illa enumeratione nasci potuit. Attractio vulgaris quadrato distantiae reciproce proportionalis, quae omnes motus coelestes tam felici successu explicat, nullius usus est nec in phaenomenis capillaribus, nec in phaenomenis adhaesionis et cohaesionis explicandis; calculus enim recte institutus facile docet, ad normam illius legis attractionem cuiusvis corporis, quocum experimenta instituere licet, i. e. cuius moles respectu totius terrae pro nihilo haberi potest, in punctum ubicunque vel adeo in contactu positum, evanescere respectu gravitatis\HistoricalNote{4}. Recte hinc concluditur, illam attractionis legem in distantiis minimis naturae haud amplius consentaneam esse, sed modificationem quandam postulare, sive quod eodem redit, corporum particulas praeter illam vim attractivam exercere aliam in distantiis minimis tantum conspicuam. Phaenomena omnia conspirant ad arguendum, hancce alteram vis attractivae partem (\textit{attractionem mole\SourceLeaf{36}cularem}), in distantiis vel minimis quas mensurare licet insensibilem esse, dum in distantiis insensibilibus partem priorem (quadrato distantiae reciproce proportionalem) longe superare possit.

Ill. \textsc{Laplace} ab hac unica suppositione circa indolem virium molecularium proficiscens, ceteroqui autem legem diminutionis pro distantiis crescentibus prorsus indeterminatam linquens, primus effectum earum in figuram superficiei fluidorum calculo accurato subiecit, et, stabilita aequatione generali pro figura aequilibrii, non modo phaenomena capillaria proprie sic dicta, sed multa alia his affinia inde explicare conatus est. Hae investigationes, per mirum cum experimentis accuratis consensum ubique confirmatae. inter pulcherrima philosophiae naturalis incrementa, quae illi magno geometrae debemus, referendae, obiectiones autem a quibusdam auctoribus contra illas directae ad maximam partem vel levis vel nullius momenti sunt\HistoricalNote{5}.

In calculis ill. \textsc{Laplace} utique occurrunt quaedam stricto argumentandi modo haud prorsus consentanea. In commentatione priori, \textit{théorie de l’action capillaire}, denotata per $\varphi f$ intensitate attractionis in distantia $f$, integrale $\int\varphi f.\mathrm{d}f$ ab $f=x$ usque ad $f=\infty$ extensum statuitur $=\Pi x$; dein integrale $\int\Pi f.f\mathrm{d}f$ ab $f=x$ usque ad $f=\infty$ extensum, $=\Psi x$; denique valores integralium $2\pi\int\Psi f.\mathrm{d}f$, $2\pi\int\Psi f.f\mathrm{d}f$ ab $f=0$ usque ad $f=\infty$ extensorum statuuntur resp. $=K$, et $=H$, denotante $\pi$ semicircumferentiam circuli pro radio $=1$. Indoles functionis $\varphi f$ prorsus intacta linquitur, dummodo insensibilis sit pro omnibus valoribus sensibilibus ipsius $f$. At ex hac \textit{sola} suppositione neutiquam sequeretur, etiam $\Pi f,\Psi f$ pro valoribus sensibilibus ipsius $f$ necessario insensibiles fieri, neque maiori iure, valores integralium $2\pi\int\Psi f.\mathrm{d}f$, $2\pi\int\Psi f.f\mathrm{d}f$ ab $f=0$ usque ad valorem sensibilem \textit{finitum} ipsius $f$ extensorum insensibiliter differre a $K,H$, uti in commentatione illa legitur; infinite multas enim formas functionis $\varphi f$ imaginari liceret, suppositioni fundamentali satisfacientes, pro quibus hae conclusiones erroneae forent. Quinadeo, si $\varphi f$ attractionem completam exprimere supponitur, revera etiam continebit partem formae $\frac{a}{ff}$, a qua attractio vulgaris pendet; sed etiamsi hic terminus pro insensibili habendus sit, dum di\SourceLeaf{37}mensiones corporum attrahentium, quales in experimentis occurrere possunt, insensibiles sunt prae tota terra, tamen iam secunda integratio, si in infinitum extenderetur, inferret functioni $\Psi f$ terminum infinitum.

At si his hisque similibus quaedam levis incuriae species subesse videtur, certe ad formam disserendi potius quam ad rem ipsam attinet. Apparet enim ex dissertatione secunda, \textit{Supplément à la théorie de l’action capillaire}, ill. \textsc{Laplace} per $\varphi f$ non attractionem completam, sed partem eam tantum, quae attractioni vulgari accedit, tacite subintellexisse; posteriorem autem nullam experimentis nostris modificationem sensibilem afferre posse, facile elucet. Quinadeo addigitat, se functionem $\varphi f$ ad instar exponentialis $e^{-if}$ considerare, denotante $i$ quantitatem permagnam, aut potius $\frac{1}{i}$ lineam perparvam. Sed ne opus quidem est, generalitatem tantopere limitare, quum is, qui rem potius quam verba intuetur, facillime videat, sufficere, si integrationes illae non in infinitum, sed tantummodo usque ad distantiam sensibilem arbitrariam, aut si mavis ad distantiam finitam dimensionibus in experimentis occurrentibus maiorem extendantur.

Alio vero defectu laborat ista theoria longe graviori, et quem quantum scimus eius cavillatores ne animadverterunt quidem. Duabus illa partibus constat. Altera stabilit aequationem generalem pro fluidi superficie libera inter differentialia partialia coordinatarum: pendet haec aequatio a vi attractiva moleculari, quam fluidi particulae in se mutuo exercent, atque haec quidem theoriae pars ita absoluta est, ut nihil essentiale desiderandum restet. Sed talis aequatio inter differentialia partialia (cuius integratio, si in analysis potestate esset, functiones arbitrarias adduceret) non sufficit ad figuram superficiei \textit{ex asse} determinandam, quod fieri nequit, nisi conditio \textit{nova} accedat indolem figurae in limitibus definiens Talem conditionem sistit pars altera theoriae, eam scilicet, ut angulus plani superficiem fluidi liberam in confiniis vasis tangentis (sive exactius, in limite vis sensibilis attractivae parietis vasis) cum plano parietem vasis ibidem tangente \textit{constans} sit, puta per relationem inter intensitates virium molecularium vasis et fluidi determinatus, siquidem continuitas figurae vasis apud confinia superficiei liberae fluidi non interrumpitur. At hancce propositionem cardinalem totius theoriae per calculum demonstrare ne suscepit quidem ill. \textsc{Laplace}; quae enim in dissertatione priori p. 5 huc spectantia afferuntur, argumentationem vagam tantummodo exhi\SourceLeaf{38}bent et quod demonstrandum erat iam supponunt: calculi autem p. 44 sq. suscepti effectu carent. In altera quidem dissertatione ascensus fluidi in tubis capillaribus per methodum aliam tractatur, cuius summa cum methodo priori collata formulam (veram utique) suppeditat pro angulo illo inter plana tangentia. Sed notare oportet, proprie hic iam \textit{supponi} quod angulus sit constans, praetereaque methodum, per se parum satisfacientem, restringi ad casum maxime specialem, ubi vas prismaticum est, parietesque verticales. His perpensis fateri oportet, theoriam ab ill. \textsc{Laplace} propositam etiamnum essentialiter mancam et incompletam esse.

Resumemus itaque ab integro theoriam figurae aequilibrii fluidorum sub actione gravitatis et virium molecularium propriarum et vasis, in quo negotio methodum prorsus diversam e primis dynamicae principiis petitam sequemur, maximamque generalitatem statim ab initio amplectemur. Haec disquisitio perducet ad insigne theorema novum, theoriam completam in unicam formulam simplicissimam contrahens, e quo utraque pars theoriae ill. \textsc{Laplace} sponte demanabit.
\par\medskip\centerline{\rule{.3\textwidth}{.3pt}}

\SourceLeaf{39}\Article{1}
Ad stabiliendam aequationem aequilibrii systematis punctorum physicorum quotcunque, quorum motus conditionibus qualibuscunque adstringuntur, maxime idoneum est principium motuum virtualium, quod sic enunciamus.

Constet systema e punctis physicis $m,m',m''$ etc., in quibus massae per easdem literas denotandae concentratae concipiantur. Sit $P$ una e viribus acceleratricibus in punctum $m$ agentibus, et dum systemati motus qualiscunque infinite parvus cum conditionibus systematis sociabilis (motus virtualis) tribui fingitur, sit $\mathrm{d}p$ motus puncti $m$ in directionem vis $P$ proiectus, i. e. per cosinum anguli, quem facit cum directione vis $P$, multiplicatus; denique sit $\Sigma P\mathrm{d}p$ aggregatum omnium similium productorum respectu omnium virium punctum $m$ sollicitantium. Perinde repraesentet $P'$ indefinite vires punctum $m'$ sollicitantes, atque $\mathrm{d}p'$ motus puncti $m'$ ad singularum directiones proiectos, similiterque de reliquis punctis. Quibus ita intellectis, conditio aequilibrii systematis consistit in eo, ut aggregatum
\[
m\Sigma P\mathrm{d}p+m'\Sigma P'\mathrm{d}p'+m''\Sigma P''\mathrm{d}p''+\ \text{etc.}
\]
pro quocunque motu virtuali fiat $=0$, uti principium motuum virtualium vulgo exprimitur, vel accuratius, in eo, ut illud aggregatum pro nullo motu virtuali adipisci possit valorem positivum.

\Article{2}
Vires hic considerandae ad tria capita reducuntur.

I. Gravitas, cuius intensitatem pro singulis punctis eandem, directiones parallelas supponere licet: illam denotabimus per $g$.

\SourceLeaf{40}II. Vires attractivae, quas puncta $m,m',m''$ etc. a se mutuo experiuntur. Intensitas attractionis functioni distantiae proportionalis sive producto huius functionis per characteristicam $f$ denotandae in massam in puncto attrahente concentratam aequalis supponitur.

III. Vires, quibus puncta $m,m',m''$ etc. ad puncta quotcunque fixa attrahuntur. Pro his viribus simili modo characteristica $F$ distantiae praefigenda utemur, et per $M,M',M''$ etc. tum puncta fixa, tum massas, quae in ipsis concentratae supponuntur, designabimus.

Quodsi iam distantiam inter bina puncta $m,m'$ per hoc signum denotamus $(m,m')$, et perinde per $(m,M)$ distantiam inter puncta $m,M$ etc., nec non per $z,z',z''$ etc. altitudines punctorum $m,m',m''$ etc. supra planum horizontale arbitrarium $H$, has partes complexus $\Sigma P\mathrm{d}p$ habebimus:
\par\begingroup\small
\[
\begin{aligned}
&-g\mathrm{d}z\\
&-m'f(m,m')\mathrm{d}(m,m')-m''f(m,m'')\mathrm{d}(m,m'')-m'''f(m,m''')\mathrm{d}(m,m''')-\ \text{etc.}\\
&-MF(m,M)\mathrm{d}(m,M)-M'F(m,M')\mathrm{d}(m,M')-M''F(m,M'')\mathrm{d}(m,M'')-\ \text{etc.}
\end{aligned}
\]
\endgroup
ubi differentialia $\mathrm{d}(m,m'),\mathrm{d}(m,m'')$ etc. sunt partialia, utpote ad solum motum virtualem puncti $m$ relatae.

Iam introducamus loco functionis $f$ eam, per cuius differentiationem oritur, puta statuatur $-fx.\mathrm{d}x=\mathrm{d}\varphi x$, sive $\int fx.\mathrm{d}x=-\varphi x$. Constans integrationis ad lubitum eligi potest; si placet (et si res fert), ita determinetur, ut fiat $\varphi\infty=0$, in quo casu $\varphi t$ exhibebit integrale $\int fx.\mathrm{d}x$ ab $x=t$ usque ad $x=\infty$ extensum. Prorsus simili modo loco functionis $F$ introducatur alia $\Phi$ talis, ut habeatur $-Fx.\mathrm{d}x=\mathrm{d}\Phi x$. Ita complexus $\Sigma P\mathrm{d}p$ fit $=$
\[
\begin{aligned}
&-g\mathrm{d}z\\
&+m'\mathrm{d}\varphi(m,m')+m''\mathrm{d}\varphi(m,m'')+m'''\mathrm{d}\varphi(m,m''')+\ \text{etc.}\\
&+M\mathrm{d}\Phi(m,M)+M'\mathrm{d}\Phi(m,M')+M''\mathrm{d}\Phi(m,M'')+\ \text{etc.}
\end{aligned}
\]
ubi notandum, differentialia in linea secunda esse partialia ad solum motum puncti $m$ relata.

At manifesto quodvis harum differentialium partialium habet supplementum suum in alio complexu. Ita tum complexus $m\Sigma P\mathrm{d}p$ tum complexus $m'\Sigma P'\mathrm{d}p'$ continet differentiale partiale $mm'\mathrm{d}\varphi(m,m')$, sed quod in priori refertur ad solum motum ipsius $m$, in posteriori ad solum motum ipsius $m'$. Hinc patet, aggre\SourceLeaf{41}gatum in art. 1 prolatum revera esse differentiale completum, et quidem $=\mathrm{d}\Omega$, si statuatur $\Omega=$
\[
\begin{aligned}
&-gmz-gm'z'-gm''z''-\ \text{etc.}\\
&+mm'\varphi(m,m')+mm''\varphi(m,m'')+mm'''\varphi(m,m''')+\ \text{etc.}\\
&\phantom{+mm'\varphi(m,m')}+m'm''\varphi(m',m'')+m'm'''\varphi(m',m''')+\ \text{etc.}\\
&\phantom{+mm'\varphi(m,m')+mm''\varphi(m,m'')}+m''m'''\varphi(m'',m''')+\ \text{etc.}\\
&\phantom{+mm'\varphi(m,m')+mm''\varphi(m,m'')+mm'''\varphi(m,m''')}+\ \text{etc.}\\
&+mM\Phi(m,M)+mM'\Phi(m,M')+mM''\Phi(m,M'')+\ \text{etc.}\\
&+m'M\Phi(m',M)+m'M'\Phi(m',M')+m'M''\Phi(m',M'')+\ \text{etc.}\\
&+m''M\Phi(m'',M)+m''M'\Phi(m'',M')+m''M''\Phi(m'',M'')+\ \text{etc.}\\
&+\ \text{etc.}
\end{aligned}
\]
Conditio aequilibrii itaque in eo consistit, ut valor functionis $\Omega$ per nullum motum virtualem accipere possit incrementum positivum, sive quod idem est, \textit{ut $\Omega$ sit maximum}.

Functionem $\Omega$ etiam sequenti modo exhibere licet:
\[
\begin{aligned}
\Omega=\Sigma m\{&-gz+\tfrac12m'\varphi(m,m')+\tfrac12m''\varphi(m,m'')+\tfrac12m'''\varphi(m,m''')+\ \text{etc.}\\
&+M\varphi(m,M)+M'\Phi(m,M')+M''\Phi(m,M'')+\ \text{etc.}\}
\end{aligned}
\]
ubi characteristica $\Sigma$ repraesentat aggregatum expressionis adscriptae cum omnibus, in quas transit, dum deinceps $m$ cum $m',m'',m'''$ etc. permutatur.

\Article{3}
Si loco punctorum discretorum $M,M',M''$ etc. assumimus corpus continuum explens spatium $S$ densitate uniformi $=C$, aggregatum
\[
M\Phi(m,M)+M'\Phi(m,M')+M''\Phi(m,M'')+\ \text{etc.}
\]
transibit in integrale $C\int\mathrm{d}S.\Phi(m,\mathrm{d}S)$ per totum spatium $S$ extendendum, denotando secundum analogiam per $(m,\mathrm{d}S)$ distantiam puncti $m$ a quovis spatii $S$ elemento $\mathrm{d}S$.

At si insuper loco punctorum discretorum $m,m',m''$ etc. corpus continuum, spatium $s$ densitate uniformi $=c$ explens, considerandum est, computus ipsius $\Omega$ integrationem duplicem requiret, atque ita perficiendus erit, ut primo pro puncto indefinito $\mu$ eruatur valor expressionis
